\section{Equivariant separating filtrations}\label{sec:filtrations}

We retain the complex $K$, the metric $\rho$, and the radii $r<R$
from Section~\ref{sec:metric}.  We will repeatedly cut $K$ into pieces of
diameter at most $2R$, lowering the dimension of the separating set at
each step.  The cuts vary over a compact probability space on which $H$
acts freely.  This allows local modifications to be separated after
taking account of the parameter, even when their translates overlap in
$|K|$.

\subsection{A free probability parameter}

The following construction is a finite-color variant of the coinduction
argument in Braun and Sauer \cite[Theorem 2.1]{BraunSauer2021}, who credit
the existence theorem to Hjorth and Molberg \cite[Theorem 0.1]{HjorthMolberg2006}.
We include the construction,
including its torsion case, so that no free-action theorem is an input.

\begin{lemma}\label{lem:free-parameter}
Every countable group $H$ admits a continuous free action on a compact
metrizable zero-dimensional space $X$, preserving a Borel probability
measure $\mu$.
\end{lemma}

\begin{proof}
For $a\in H\setminus\{1\}$, put $k_a=\operatorname{ord}(a)$ if $a$ has
finite order, and $k_a=2$ otherwise.  Let
\[
 X_a=\{\xi:H\longrightarrow\Z/k_a\Z:
                 \xi(ta)=\xi(t)+1\text{ for every }t\in H\}.
\]
This is a closed subspace of the product of finite discrete spaces
$(\Z/k_a\Z)^H$.  Choose a representative $t_C$ for each left coset
$C=t_C\langle a\rangle$.  A coloring in $X_a$ is specified uniquely by
offsets $\omega_C\in\Z/k_a\Z$, through
$\xi(t_Ca^j)=\omega_C+j$.  This is well defined also when $a$ has finite
order, because the modulus is then its order.  Thus $X_a$ is nonempty,
compact and zero-dimensional.

Give the offsets independent uniform distributions.  The left action
\[
 (h\xi)(t)=\xi(h^{-1}t)
\]
preserves $X_a$ and its probability measure: it permutes the cosets and
adds a fixed element of $\Z/k_a\Z$ to each offset.  These operations
preserve every finite cylinder probability.  Take the product over
$a\ne1$ to obtain $X$ and $\mu$.  Countability gives metrizability, and
the coordinatewise action is continuous and measure preserving.
For any $a\ne1$ and any point of $X$, its $a$-coordinate satisfies
\[
 (a\xi_a)(1)=\xi_a(a^{-1})=\xi_a(1)-1\ne\xi_a(1).
\]
Hence no nonidentity element fixes a point.  If $H$ is trivial, take a
one-point probability space.
\end{proof}

Fix such $(X,\mu)$ for the rest of the geometric argument.  In
particular, no amenability assumption on $H$ is needed.

\subsection{The allowed filtrations}

A triangulation refining $K$ means a finite affine simplicial
subdivision of each closed original simplex, with agreement on shared
faces.  A family $T(x)$ of such triangulations, together with specified
subcomplexes, is \emph{locally constant} if, on each original simplex,
all its exact affine simplices and subcomplex markings are constant on
the pieces of a finite clopen partition of $X$.  The partition may
depend on the original simplex.  We require equivariance:
$T(hx)=hT(x)$, and likewise for every specified subcomplex.
Refinement of triangulations always preserves their earlier
subcomplexes.

Begin with $T_n(x)=K$ and $Z_n(x)=|K|$.  At stage $d$, where
$n\ge d\ge1$, suppose that $P(x)=Z_d(x)$ is a subcomplex of $T_d(x)$.
An \emph{admissible} next choice consists of a locally constant equivariant
refinement $T_{d-1}(x)$ and its subcomplex $Z_{d-1}(x)$ such that
\begin{enumerate}[label=\textup{(\roman*)}]
\item $Z_{d-1}(x)\subseteq P(x)$;
\item for every simplex $F$ of $T_d(x)$ contained in $P(x)$,
\[
 \dim\bigl(F\cap Z_{d-1}(x)\bigr)\le\dim F-1;
\]
\item every connected component of $P(x)\setminus Z_{d-1}(x)$ has
$\rho$-diameter at most $2R$.
\end{enumerate}
A negative dimension bound means that the intersection is empty.
Thus the new separator avoids the vertices of $T_d$ and meets each
higher-dimensional face in lower dimension; its own vertices belong to
the finer triangulation $T_{d-1}$.
Connected components use the subspace topology given by $\rho$.
We call a filtration satisfying these conditions admissible.
In particular, these conditions imply $\dim Z_i(x)\le i$.

The local operation behind the construction is easiest to describe at a
fixed parameter. For a separator $Z\subseteq P$, a center $c\in|K|$,
and $t\in(r,R)$, remove the part of $Z$ inside the ball and insert
the corresponding slice of $P$:
\[
 Z^{\prime}=(Z\setminus U_t(c))\cup(P\cap S_t(c)).
\]
The inserted slice prevents a component of $P\setminus Z^{\prime}$
from passing between the ball and its exterior. Inside the ball its
diameter is at most $2t$; outside, the old separator remains.
The next lemma makes this operation compatible with triangulations,
group translates, and clopen parameter choices. It also constructs the
level cuts used to start a filtration. Only the affine descriptions of
distance levels enter the construction; no Euclidean metrics need to
agree across faces.

\begin{lemma}[Compatible cuts]\label{lem:refinement}
Fix locally constant equivariant triangulation families $T$ and $S$,
where $S$ refines $T$, and subcomplex families $P\subseteq T$ and
$Z\subseteq S$ with $Z\subseteq P$.  Let
$\{(c_\alpha,B_\alpha)\}_{\alpha\in A}$ be a countable equivariant
family of centers $c_\alpha\in|K|$ and clopen sets $B_\alpha\subseteq X$.
Assume that only finitely many of the balls $L_R(c_\alpha)$ meet any
given original simplex.  Set
\[
 \mathcal U_t(x)=\bigcup_{\alpha:\,x\in B_\alpha}U_t(c_\alpha),
 \qquad
 \mathcal S_t(x)=\bigcup_{\alpha:\,x\in B_\alpha}S_t(c_\alpha).
\]
Outside a countable set of radii $t\in(r,R)$, there is a locally
constant equivariant triangulation refining $S$ in which both
\begin{equation}\label{eq:cut-replacement}
\begin{gathered}
 P(x)\cap\mathcal S_t(x),\\
 \bigl(Z(x)\setminus\mathcal U_t(x)\bigr)
       \cup\bigl(P(x)\cap\mathcal S_t(x)\bigr)
\end{gathered}
\end{equation}
are subcomplexes.  Each inserted level meets every face $F$ of $T$ or
$S$ in dimension at most $\dim F-1$.  Consequently, if $Z$ satisfies
the dimension-drop condition relative to $T$, both displayed sets do
also.  If every component of $P\setminus Z$ has diameter at most $2R$,
the second displayed set preserves that property.
\end{lemma}

\begin{proof}
On a fixed original simplex the function
$f_c(p)=\rho(c,p)$ is a maximum of finitely many affine functions.
For a fixed face $F$, a level $\{f_c=t\}$ can have full dimension in
$F$ only if one of these affine functions is identically $t$ on the
affine span of $F$.  Otherwise the level is contained in a finite
union of proper affine hyperplanes.  Thus only finitely many levels
are exceptional for this face and center.  There are countably many
original simplices and centers, and only finitely many exact
triangulation types on each original simplex.  Excluding the
resulting countable union of exceptional levels gives the stated
dimension drop simultaneously for all faces of $T$ and $S$.

We give the compatible refinement explicitly.  On each original
simplex $\sigma$, cut the simplices of $S$ by all hyperplanes $b=t$,
where $b$ is an affine branch for one of the finitely many relevant
centers.  Use the
canonical branch collection from the finite potential polytope in
Section~\ref{sec:metric}.  On each piece of the finite clopen partition
recording the old data and active conditions, this gives a finite
polytopal complex $\mathcal P_\sigma$.  It represents the old
subcomplexes and the closed conditions defining the two sets in
\eqref{eq:cut-replacement}.

These polytopal complexes may initially disagree on original faces.
For each original simplex $\tau$, let $\mathcal Q_\tau$ be the common
polytopal refinement of
\[
 \mathcal P_\sigma|_\tau
 \quad\text{over all original cofaces }\sigma\supseteq\tau,
 \quad\text{including }\sigma=\tau.
\]
There are finitely many cofaces by Lemma~\ref{lem:proper}.
The common refinement consists of the nonempty intersections of cells
and their faces.  If $\tau$ is a face of $\sigma$, then
$\mathcal Q_\tau$ refines $\mathcal Q_\sigma|_\tau$, since every coface
of $\sigma$ is also a coface of $\tau$.

Triangulate the original simplices in increasing dimension.  When
handling $\sigma$, its boundary has already been triangulated by a
refinement of the boundary restrictions of $\mathcal Q_\sigma$.
Extend this triangulation over the cells of $\mathcal Q_\sigma$ in
increasing cell dimension.  Keep cells in the original boundary as
already triangulated.  For every other positive-dimensional cell,
cone its triangulated boundary to the arithmetic mean of its vertices,
which lies in its relative interior.  Zero-dimensional cells need no
operation.  The cones agree on common faces and give the desired finite
affine subdivision, preserving the prescribed boundary triangulation.

Every operation commutes with the affine maps induced by $H$:
restriction, common refinement, and the mean of the full vertex set of
a polytope are all canonical.  This proves equivariance even for
nontrivial simplex stabilizers.  On each original simplex the procedure
uses only finitely many data: those on its faces and their original
cofaces, together with finitely many relevant cuts.  A common finite
clopen partition makes all these data constant.  Thus the output has
the required local constancy.  No manifold condition on $K$ is involved.

The first set in \eqref{eq:cut-replacement} has the dimension bound
already proved for the levels.  The second is contained in the union
of those levels with $Z$, so it has the same bound whenever $Z$ does.
For its separation property, let $C$ be a component of its complement
in $P(x)$.  If $C$ meets an active $U_t(c_\alpha)$, continuity of
$f_{c_\alpha}$ and connectedness of $C$ imply that $C$ stays in that
ball, since it avoids the inserted level.  Its diameter is then at
most $2t<2R$.  Otherwise $C$ lies outside all active open balls, where
the old set $Z$ was retained.  It is therefore contained in an old
component of $P(x)\setminus Z(x)$ and has diameter at most $2R$.
\end{proof}

The genericity against $S$ is useful when $S$ already contains the old
separator $Z$: intersections of an inserted level with its faces have
one less dimension.  The later area comparison will use precisely
this conclusion.  The lemma itself does not require disjoint active
balls; disjointness will be imposed when their areas are counted.

\begin{lemma}\label{lem:filtration-exists}
For any allowed choices through stage $d$, an allowed choice of
$(T_{d-1},Z_{d-1})$ exists.  Hence a filtration can be continued to
$Z_0$.
\end{lemma}

\begin{proof}
Each original simplex is compact in $\rho$.  Cover representatives of
the finitely many original simplex orbits by finitely many open balls
$U_r(c)$, and take all their $H$-translates.  These balls cover $|K|$.
By Lemma~\ref{lem:proper}, only finitely many translated radius-$R$
balls meet a given original simplex.  Apply Lemma~\ref{lem:refinement}
with every activation set equal to $X$, with $T=S=T_d$, and with
$Z=\varnothing$.  For a common allowed radius $t\in(r,R)$, take
\[
 Z_{d-1}(x)=Z_d(x)\cap\bigcup_c S_t(c).
\]
This is a subcomplex of the resulting refinement and has the required
dimension drop.  If $C$ is a component of $Z_d(x)\setminus Z_{d-1}(x)$,
choose $p\in C$ and a covering ball $U_r(c)$ containing it.
The connected set $f_c(C)$ avoids $t$ and contains a value below $t$,
so $C\subseteq U_t(c)$.  Its diameter is at most $2t<2R$.
The assertion is vacuous for an empty component space.
\end{proof}

\subsection{Dimension drop along the entire filtration}

We record the consequence that will control chains in
Section~\ref{sec:chains}.  For any simplex $F$ of $T_i(x)$ contained in
$Z_i(x)$ and any $0\le j\le i$,
\begin{equation}\label{eq:iterated-drop}
 \dim\bigl(F\cap Z_j(x)\bigr)\le\dim F-(i-j).
\end{equation}
Indeed, the first step is the defining condition.  Its intersection is
a subcomplex in the next refinement.  Applying the next dimension
drop to each simplex of that subcomplex decreases its maximal
dimension by one.  Induction gives \eqref{eq:iterated-drop}.
Applied with $i=n$, this also shows that a proper original face meets
$Z_j(x)$ in dimension strictly less than $j$.

\begin{lemma}[Carriers of the zero-dimensional stratum]\label{lem:carrier}
Let $(T_i,Z_i)_{i=0}^n$ be an allowed filtration and let $z\in Z_0(x)$.
For every $i$, there is a unique $i$-simplex of $T_i(x)$ contained in
$Z_i(x)$ and containing $z$.  The point $z$ lies in its relative
interior.  In particular, every point of $Z_0(x)$ lies in the interior
of an original top simplex.
\end{lemma}

\begin{proof}
Let $F$ be the unique simplex of $T_i(x)$ whose relative interior
contains $z$.  Since $Z_i(x)$ is a subcomplex containing $z$, it contains
$F$, and $\dim F\le i$.  If $\dim F<i$, then
\eqref{eq:iterated-drop} with $j=0$ would make $F\cap Z_0(x)$ empty.
Thus $\dim F=i$.  Every other simplex containing $z$ has $F$ as a face,
so no other $i$-simplex contains it.  Taking $i=n$ proves the last
assertion.
\end{proof}

We have obtained a nonempty class of equivariant separating
filtrations, together with a local cutting operation that stays within
that class.  The next section chooses separators with area close to
the infimum and compares them with these cuts.
